\chapter*{Введение}
\addcontentsline{toc}{chapter}{Введение}
\label{ch:intro}

Дифференцируемый рендеринг, то есть вычисление производных изображения по параметрам 3D-сцены, позволяет решать обратные задачи компьютерной графики и зрения градиентной оптимизацией: по фотографиям восстанавливаются геометрия объектов, их материалы и источники освещения. Существующие методы этого класса можно условно разделить на две группы. Методы на основе нейронных и объёмных представлений (NeRF, 3D Gaussian Splatting и их развития) быстро восстанавливают внешний вид сцены, но не дают явной поверхности и физически интерпретируемых параметров освещения. Физически обоснованные методы дифференцируемого рендеринга функций расстояния со знаком и систем общего назначения (Mitsuba~3) восстанавливают физические параметры сцены, однако требуют больших вычислительных затрат: плотные сетки функций расстояния растут в 8 раз при каждом удвоении разрешения и уже на разрешении $512^3$ не помещаются в память распространённых графических процессоров, а системы общего назначения заново трассируют и дифференцируют световые пути на каждом шаге оптимизации, даже когда геометрия сцены известна. Подробные обзоры приведены в разделах~\ref{ch:chap1} и~\ref{ch:light-review}.

Реконструкция геометрии и освещения по фотографиям востребована в инженерном проектировании, в проектировании систем освещения, в кинематографе и в построении цифровых двойников. Высокие требования существующих методов к памяти и времени вычислений ограничивают точность реконструкции доступным разрешением и не позволяют использовать их на массовом оборудовании. Задача реконструкции источника освещения по фотографии с учётом многократных отражений света в сцене и получения результата в стандартном формате фотометрических данных (IES) непосредственно связана с практическими задачами проектирования освещения.

Предложен метод дифференцируемого рендеринга поверхностей, использующий для физически обоснованной реконструкции разреженное множество бриков функции расстояния со знаком вместе с адаптированными к нему алгоритмами поиска точек релаксированной границы и редистансинга. Предложен метод реконструкции источников освещения с гониометрической диаграммой направленности и пространственно неоднородных площадных источников, который использует линейность изображения по параметрам излучения при известной геометрии: световые пути трассируются один раз, после чего изображение и его градиент вычисляются проходами по кэшу световых путей, а при фиксированной угловой модели задача сводится к линейной задаче наименьших квадратов.

Работа продолжает первый этап проекта, на котором были разработаны компактные представления поверхностей, материалов и источников освещения. Разреженное множество бриков, используемое в методе дифференцируемого рендеринга поверхностей, является одним из таких компактных представлений, а разработанный метод позволяет получать компактные модели реконструкцией по изображениям. Метод реконструкции освещения использует систему расчёта освещённости Hydra и транслятор kernel\_slicer, порождающий код для графических процессоров на языке Slang для Vulkan, что обеспечивает кроссплатформенность реализации.

Проект выполняется в два этапа: 2025 год (первый этап) и 2026 год (второй, заключительный этап). Общая цель проекта -- разработка новых методов и алгоритмов прямого и обратного рендеринга, достигающих многократного выигрыша по вычислительным затратам и затратам по памяти у существующих методов, а также обеспечение их доступности на широком спектре вычислительных устройств за счёт создания кроссплатформенного программного обеспечения на базе собственных методов и алгоритмов. Цель второго этапа -- разработка эффективных методов дифференцируемого рендеринга (задача 2.2.2 ТЗ), позволяющих восстанавливать по изображениям поверхности и конфигурацию освещения. В соответствии с п.~5.2.1 ТЗ на этапе решаются следующие задачи:
\begin{itemize}
    \item проведение аналитического обзора научной литературы на тему эффективных методов дифференцируемого рендеринга для реконструкции или оптимизации конфигурации освещения (раздел~\ref{ch:light-review}) и дифференцируемого рендеринга поверхностей (раздел~\ref{ch:chap1});
    \item разработка и реализация метода дифференцируемого рендеринга для реконструкции или оптимизации конфигурации освещения (раздел~\ref{ch:light-method}) и проведение его сравнения с аналогами (раздел~\ref{chap:light-comparison});
    \item разработка и реализация метода дифференцируемого рендеринга поверхностей (раздел~\ref{ch:chap2}) и проведение его сравнения с аналогами (раздел~\ref{ch:chap4});
    \item разработка комплекса методов и алгоритмов, объединённых в систему LightCadRenderCore, для рендеринга и реконструкции освещения (подраздел~\ref{sec:lightcad}).
\end{itemize}
Результаты этапа завершают проект: компактные представления первого этапа дополняются методами, позволяющими получать их реконструкцией по изображениям, а методы реконструкции освещения -- системой, применимой в задачах проектирования освещения.

\endinput
